\documentclass[11pt]{amsart}
\usepackage{calc,amssymb,amsmath,ulem,stmaryrd,fullpage}
\usepackage{enumerate}
\usepackage{alltt}
\RequirePackage[dvipsnames,usenames]{color}
\let\mffont=\sf
\normalem
%\input{kmacros3.sty}
\input{xy}
\xyoption{all}
\usepackage{hyperref}
\usepackage{graphicx}
\usepackage[all,cmtip]{xy}
\usepackage{verbatim}
\usepackage[left=0.95in,top=0.95in,right=0.95in,bottom=0.95in]{geometry}
\newcommand{\tauCohomology}{T}
\newcommand{\FCohomology}{S}


\theoremstyle{theorem}
%\newtheorem*{mainthm*}{Main Theorem}

\newcommand{\note}[1]{\marginpar{\sffamily\tiny #1}}

\def\todo#1{\textcolor{Mahogany}%
{\sffamily\footnotesize\newline{\color{Mahogany}\fbox{\parbox{\textwidth-15pt}{#1}}}\newline}}

\renewcommand{\O}{\mathcal O}
\newcommand{\ie}{{\itshape i.e.} }
\newcommand{\roundup}[1]{\lceil #1 \rceil}
\newcommand{\rounddown}[1]{\lfloor #1 \rfloor}
\renewcommand{\to}[1][]{\xrightarrow{\ #1\ }}
\newcommand{\onto}[1][]{\protect{\xrightarrow{\ #1\ }\hspace{-0.8em}\rightarrow}}
\newcommand{\into}[1][]{\lhook \joinrel \xrightarrow{\ #1\ }}
%\newcommand{DBDual}[1]{{\underline{\omega}_{#1}}}
\newcommand{\DBDual}[1]{{{\uuline \omega}{}^{\mydot}_{#1}}}
\newcommand{\Tt}{{\mathfrak{T}}}
%\newcommand{\bn}{{\mathfrak{n}} }
\newcommand{\sol}[1]{ \vskip 6pt{\color{blue} {\bf Solution:} #1} \vskip 6pt}

\newcommand{\bR}{{\mathbb{R}}}
\newcommand{\va}{{\vec{a}}}
\newcommand{\vb}{{\vec{b}}}
\newcommand{\vzero}{{\vec{0}}}
\newcommand{\vi}{{\vec{i}}}
\newcommand{\vj}{{\vec{j}}}
\newcommand{\vk}{{\vec{k}}}
\newcommand{\vh}{{\vec{h}}}
\newcommand{\vr}{{\vec{r}}}
\newcommand{\vu}{{\vec{u}}}
\newcommand{\vv}{{\vec{v}}}
\newcommand{\vw}{{\vec{w}}}
\newcommand{\vx}{{\vec{x}}}
\newcommand{\vy}{{\vec{y}}}
\newcommand{\vz}{{\vec{z}}}
\newcommand{\vT}{{\vec{T}}}
\newcommand{\vN}{{\vec{N}}}
\newcommand{\vF}{{\vec{F}}}
\newcommand{\vG}{{\vec{G}}}
\newcommand{\bF}{\mathbb{F}}
\newcommand{\bQ}{\mathbb{Q}}
\newcommand{\bZ}{\mathbb{Z}}
\newcommand{\bC}{\mathbb{C}}
\newcommand{\Gal}{\textnormal{Gal}}
\newcommand{\Aut}{\textnormal{Aut}}
\newcommand{\Emb}{\textnormal{Emb}}
\newcommand{\Tr}{\textnormal{Tr}}
\newcommand{\N}{\textnormal{N}}
\newcommand{\Hom}{\textnormal{Hom}}
\newcommand{\Ext}{\textnormal{Ext}}
\newcommand{\Spec}{\textnormal{Spec}}
\newcommand{\fram}{\mathfrak{m}}
\newcommand{\id}{\textnormal{id}}


\begin{document}
\title{HW \#7 -- Math 6320\\Spring 2015}
\author{Due: Thursday April 16th}
\maketitle

\begin{enumerate}
\item Suppose that $M$ is a finitely generated module over a $R$.  Suppose that $f : M \to M$ is a surjective $R$-module homomorphism.  Show that $f$ is an isomorphism.
    \vskip 3pt
    \emph{Hint: } Nakayama's lemma is a good approach.  View $M$ as an $R[x]$ module where $x$ acts via $f$.
    \vskip 6pt
\item Suppose that $M$ is a finitely generated projective module over a local ring $(R, \mathfrak{m})$.  Show that $M$ is actually a free module.  Another way to say this is that \emph{projective modules are locally free}.  
    \vskip 3pt
    \emph{Hint: } Use the previous problem.
    \vskip 6pt
\item Suppose that $R$ is a ring and $M$ is an $R$-module.  Show that $M$ is injective if and only if for every ideal $J \subseteq R$ and every $R$-module homomorphism $g : J \to M$, $g$ can be extended to $g' : R \to M$.  
    \vskip 3pt
    \emph{Hint: } Replacing $J \subseteq R$ with an arbitrary inclusion of modules gives you the definition of \emph{injective}.  Hence say $N' \subseteq N$ and $g : N' \to M$ is a map.  Consider the set of submodules of $N$ to which $g$ extends, use Zorn's lemma to make a biggest submodule of $N$ to which $g$ extends and then make a bigger one for a contradiction.
    \vskip 6pt
\item Suppose that $R$ is a ring.
\begin{itemize}
\item[(a)]  If $W\subseteq R$ is a multiplicative set and $I$ is an injective $R$-module, show that $W^{-1} I$ is an injective $W^{-1} R$-module.
\item[(b)]  Show that an arbitrary direct sum $\bigoplus_\lambda I_{\lambda}$ of injective $R$-modules is an injective $R$-module.
\end{itemize}
\vskip 6pt
\item Suppose that $\phi : R \to S$ is a ring homomorphism and $I$ is an injective $R$-module.  Show that $\Hom_R(S, I)$ is an injective $S$-module.
    \vskip 3pt
    \emph{Hint: } This is not so hard, you need to prove/use a special case of a jazzed up version of $\Hom-\otimes$ adjointness:  for any $S$-module $M$, $\Hom_R(M, I) \cong \Hom_S(M, \Hom_R(S, I))$ as $S$-modules.  
\vskip 6pt
\item If $R$ is a PID.  Show that an $R$-module $I$ is injective if and only if it is divisible\footnote{For an integral domain $R$, an $R$-module $M$ is called divisible if for every $0 \neq r \in R$, and every $m \in M$, there is some $m' \in M$ with $rm' = m$.}.
    \vskip 6pt
\item Show that every $\bZ$-module can be embedded into an injective $\bZ$-module
\vskip 3pt
\emph{Hint: } First show that a free $\bZ$-module can be embedded into an injective $\bZ$-module.  For $M$ an arbitrary $\bZ$-module, write $F \to M \to 0$ where $F$ is a free $\bZ$-module and $F \to M$ has kernel $K$.  Mod out some appropriate injective module by $K$.
\vskip 6pt
\item Conclude that for any ring $R$, any $R$-module can be embedded into an injective $R$-module.  Hence show that injective resolutions exist.  This result is an assertion that \emph{the category of $R$-modules has enough injectives}.
\end{enumerate}
Continued on the next page.
\begin{enumerate}
\setcounter{enumi}{9}
\item Suppose that $M, N$ are $R$-modules and that $\Ext^1(M, N) = 0$.  Show that every extension $0 \to N \to E \to M \to 0$ is split.  More generally, identify a particular (obstruction) class in $\Ext^1(M, N)$ which is zero if and only if a given extension $0 \to N \to E \to M \to 0$ is split.
    \vskip 3pt
    \emph{Hint:} Compute the long exact sequence associated to $\Hom(M, \bullet)$ and think about what it means.   The class should be the image of particular identity map in $\Ext^1$.  
\vskip 6pt
\item We will show that $\Ext^1(M, N)$ is in bijection with the set of equivalence classes of extensions $0 \to N \to E \to M \to 0$.\footnote{Remember, two extensions are equivalent if there is a commutative diagram 
    \[
        \xymatrix{
            0 \ar[r] & N \ar[d]_{\id_N} \ar[r] & E \ar[d]_{\sim} \ar[r] & M \ar[r] \ar[d]_{\id_M}  & 0 \\
            0 \ar[r] & N \ar[r] & E' \ar[r] & M \ar[r] & 0 
            }
    \]}
    \begin{itemize}
    \item[(a)]  Start with an exact sequence $0 \to K \xrightarrow{i} P \xrightarrow{\rho} M \to 0$ with $P$ projective and induce a map $\partial : \Hom(K, N) \to \Ext^1(M, N) \to 0$.  Thus each class $x$ in $\Ext^1$ gives us a (non-canonical) $\beta \in \Hom(K, N)$ with $\partial(\beta) = x$.  Let $X$ be the cokernel of $K \to P \oplus N$ defined by $k \mapsto (i(k), \beta(k))$.  Prove that there is a commutative diagram
        \[
        \xymatrix{
            0 \ar[r] & K \ar[r] & P \ar[d] \ar[r] & M \ar[d]_{\id_M} \ar[r] & 0 \\
            0 \ar[r] & N \ar[r] & X \ar[r] & M \ar[r] & 0 
        }
        \]
        with exact rows.  Explain in particular what the map $X \to M$ is.  
    \item[(b)]  In the previous problem, you constructed an element in $\Ext^1(M, N)$ which was zero if and only if a given extension was split.  Show that the element $x \in \Ext^1$ corresponds to the extension $0 \to N \to X \to M \to 0$.  
    \item[(c)]  Use what you have already done to complete the proof of the desired statement.  
    \end{itemize}
\end{enumerate}
\end{document}

